\documentclass{article}



% Language setting

% Replace `english' with e.g. `spanish' to change the document language

\usepackage[english]{babel}



% Set page size and margins

% Replace `letterpaper' with `a4paper' for UK/EU standard size

\usepackage[letterpaper,top=2cm,bottom=2cm,left=3cm,right=3cm,marginparwidth=1.75cm]{geometry}



% Useful packages

\usepackage{amsmath}

\usepackage{amssymb}

\usepackage{graphicx}

\usepackage[colorlinks=true, allcolors=black]{hyperref}

\usepackage{titlesec}
\usepackage{xparse}



\NewDocumentCommand{\qcq}{m o o m}{%
	\ensuremath{\forall #1 \IfValueT{#2}{, #2} \IfValueT{#3}{, #3}  \in  \mathbb{#4}}%
}
\NewDocumentCommand{\ext}{m o o m}{%
	\ensuremath{\exists #1 \IfValueT{#2}{, #2} \IfValueT{#3}{, #3}  \in  \mathbb{#4}}%
}



\title{\textbf{Exercices sur les Limites}}

\author{\textbf{Yassin Akkab}}



\newcommand{\R}{\mathbb{R}}
\newcommand{\N}{\mathbb{N}}
\newcommand{\Z}{\mathbb{Z}}
\newcommand{\C}{\mathbb{C}}
\newcommand{\Q}{\mathbb{Q}}
\newcommand{\D}{\mathbb{D}}





\graphicspath{ {./Screenshots/} }




\begin{document}
	
	\maketitle
	
	

	
	
	\subsection*{Exercice 1}
	Calculez les limites suivantes :
	\[
	\lim_{x \to 0} \frac{\sin(x)}{x}
	\]
	\[
	\lim_{x \to +\infty} \frac{1}{x^2}
	\]
	\[
	\lim_{x \to 1} \frac{x^2 - 1}{x - 1}
	\]
	
	\subsection*{Exercice 2}
	Étudiez les limites suivantes, en précisant s'il s'agit de formes indéterminées. Si oui, appliquez la règle de L'Hôpital si nécessaire.
	\[
	\lim_{x \to +\infty} \frac{e^x}{x^n} \quad \text{pour tout} \quad n \in \mathbb{N}
	\]
	\[
	\lim_{x \to 0^+} \frac{1}{x}
	\]
	\[
	\lim_{x \to 0^+} \frac{\sin(x)}{x}
	\]
	
	\subsection*{Exercice 3}
	Montrez que :
	\[
	\lim_{x \to 0^+} \frac{1}{x} = +\infty
	\]
	\[
	\lim_{x \to 0^+} \frac{x^2}{x} = 0
	\]
	
	\subsection*{Exercice 4}
	Calculez :
	\[
	\lim_{x \to 0^+} \frac{e^x - 1}{x}
	\]
	\[
	\lim_{x \to +\infty} \frac{x}{e^x}
	\]
	
	\subsection*{Exercice 5}
	Comparez les comportements asymptotiques des fonctions suivantes :
	\[
	\lim_{x \to +\infty} \frac{x^n}{e^x} \quad \text{pour tout} \quad n \in \mathbb{N}
	\]
	\[
	\lim_{x \to +\infty} \frac{x^n}{x^{n+1}}
	\]
	
		
		
	
	
	
	
	
	
	
\end{document}